% ATTEMPT_STATUS: unresolved
% TOP_PROBLEM: {"problemId": "problem.schinzel-s-hypothesis-h-on-simultaneous-prime-values-of-irreducible-polynomials", "canonicalTitle": "Schinzel's Hypothesis H on Simultaneous Prime Values of Irreducible Polynomials", "releaseRank": 54, "primaryDomain": "number_theory_arithmetic_geometry", "primaryDomainLabel": "Number theory & arithmetic geometry", "displayStatus": "open_with_solved_subcases", "releaseStatus": "open_with_solved_subcases"}
% !TeX program = lualatex
% Complete problem section copied verbatim from research_batch_51_54.tex.
\documentclass[11pt]{article}
\usepackage[margin=25mm]{geometry}
\usepackage{fontspec}
\setmainfont{Latin Modern Roman}
\usepackage{amsmath,amssymb,mathtools}
\usepackage{unicode-math}
\setmathfont{Latin Modern Math}
\usepackage{enumitem,needspace}
\usepackage{xurl,hyperref}
\hypersetup{colorlinks=true,urlcolor=blue}
\setcounter{secnumdepth}{1}
\setlength{\parindent}{0pt}
\setlength{\parskip}{5pt}
\setlength{\emergencystretch}{3em}
\setlist[itemize]{leftmargin=3em,itemsep=5pt,topsep=4pt,parsep=0pt}
\allowdisplaybreaks
\newcommand{\N}{\mathbb{N}}
\newcommand{\Z}{\mathbb{Z}}
\newcommand{\Q}{\mathbb{Q}}
\newcommand{\R}{\mathbb{R}}
\newcommand{\C}{\mathbb{C}}
\newcommand{\F}{\mathbb{F}}
\newcommand{\A}{\mathbb{A}}
\newcommand{\PP}{\mathbb P}
\newcommand{\E}{\mathbb{E}}
\newcommand{\eps}{\varepsilon}
\newcommand{\dd}{\,\mathrm d}
\newcommand{\sref}[2]{\hyperlink{source:#1:#2}{[S#2]}}
\newcommand{\eref}[2]{\hyperlink{extra:#1:#2}{[E#2]}}
% Turn the next switch off for a shorter reading copy. The quotations preserve
% every exactTarget field, including qualifications omitted from a short summary.
\newif\ifshowcatalogue
\showcataloguetrue
\newcommand{\cataloguescope}[1]{\ifshowcatalogue
 \paragraph{Exact scope in the supplied catalogue.}
 \begin{quote}\small #1\end{quote}\fi}
\begin{document}
\setcounter{section}{53}
% ================================================================
% problemId: problem.schinzel-s-hypothesis-h-on-simultaneous-prime-values-of-irreducible-polynomials
% source releaseRank: 54; source releaseStatus: open_with_solved_subcases
% exactTarget SHA-256: f079951f6ce8970223c01331fea5a722075ed301c3a6eacabd9032382ca38769
\clearpage
\section{\texorpdfstring{Schinzel's Hypothesis H on Simultaneous Prime Values of Irreducible Polynomials}{Schinzel's Hypothesis H on Simultaneous Prime Values of Irreducible Polynomials}}\label{54}
\subsection{Definitions and mathematical statement}
Let $f_1,\ldots,f_r\in\Z[x]$ be nonconstant irreducible polynomials with positive leading coefficients. The local admissibility condition is
\[
 \forall p\text{ prime}\ \exists n\in\Z:\quad p\nmid\prod_{i=1}^r f_i(n).
\]
The proposed conclusion is that $\{n\ge1:f_i(n)\text{ is prime for every }i\}$ is infinite. Irreducibility is over $\Q$ for primitive integer polynomials; local admissibility excludes fixed prime divisors. No asymptotic frequency is included in this qualitative assertion.
\par\smallskip\textit{Formulation and concept references:} \sref{54}{1}.
\cataloguescope{Let f1,...,fr in Z[x] be irreducible with positive leading coefficients and no prime dividing f1(n)...fr(n) for every n. Prove that infinitely many positive integers n make every fi(n) prime.}
\subsection{Short English statement}
If irreducible integer polynomials have no fixed local obstruction, must they simultaneously take prime values infinitely often?
% BEGIN RESEARCH ADDITION ATTEMPT-54
\subsection{Research attempt}

\paragraph{Current frontier.}
For a single linear polynomial the conclusion is Dirichlet's theorem;
for the general fixed family the available coefficient-averaging results
have different quantifiers. Skorobogatov--Sofos, Theorems~1.1--1.2,
show that almost all admissible tuples, ordered by coefficient height,
have at least one simultaneous prime value. Their introductory warning
that an almost-all one-value assertion cannot simply be iterated into a
fixed-polynomial infinitude assertion is important here. \eref{54}{1}
Browning--Sofos--Ter\"av\"ainen, Theorems~1.2 and~1.4 in the May 2026
revision, obtain averaged prime-value asymptotics with arbitrary fixed
logarithmic power savings in the exceptional proportion. Their averaging
is over coefficient boxes, not the affine substitution orbit of an
arbitrarily prescribed tuple. \eref{54}{2}

The 2026 Schinzel--W\'ojcik paper assumes Hypothesis H; it is not a proof
of it. \sref{54}{2}
The catalogue's Golomb-method preprint claims a solution, but a specific
identity and a limiting inference fail the checks below; neither is used
as an established result. \sref{54}{3}, \eref{54}{4}
The Scavia download in [S1] could not be retrieved in this run, so its
bibliographic details were not independently verified. The argument uses
the complete explicit fixed-family formulation printed in this section,
which agrees with the formulation inspected in the introduction of
\eref{54}{1}; no uninspected source-only hypothesis is added.

\paragraph{Attack route.}
A finite covering by prime divisors cannot give a counterexample: local
admissibility and the Chinese remainder theorem always leave a residue
class avoiding any prescribed finite set of primes. Sieve methods would
need additional parity-sensitive information to turn such avoidance
into actual primes; the asymptotic sieve makes explicit that additional
analytic input, rather than ordinary congruence counts alone, can matter.
\eref{54}{3}

The route pursued here tries to transfer a coefficient-average theorem
to a fixed tuple without making the quantifier error just described.
A single counterexample is stable under many substitutions $x\mapsto
ax+b$. If sufficiently many of those substitutions retain admissibility,
a counterexample forces a quantitatively large exceptional set in
coefficient space. The questions are then exact: how many admissible
substitutions exist, how much do they increase height, and is the forced
exceptional set larger than the available upper bounds? The first two
questions are answered below; the third exposes the unresolved gap.

\paragraph{Attempt: density of admissible affine substitutions.}
Remove exact repeated polynomials, which does not change simultaneous
primality. Keep the remaining ordered tuple
$\mathbf f=(f_1,\ldots,f_r)$ fixed, with the hypotheses in the original
statement. Put
\[
 \begin{gathered}
 F=\prod_{i=1}^r f_i,\qquad
 \Delta=\deg F=\sum_i d_i,\qquad D=\max_i d_i,\\
 \nu_F(p)=\#\{b\bmod p:F(b)\equiv0\pmod p\}.
 \end{gathered}
\]
Admissibility gives $\nu_F(p)<p$ and ensures that $F$ is not the zero
polynomial modulo any prime; hence $\nu_F(p)\le\Delta$ as well.

\textbf{Lemma 54.1 (affine parameter density).} Define
\[
 N_F(T)=\#\{1\le a,b\le T:\gcd(a,F(b))=1\}.
\]
Here $T$ is a positive integer and gcd is taken with the absolute value
when needed. Then
\begin{equation}\label{ra:54:density}
 N_F(T)=\bigl(c_F+o(1)\bigr)T^2,
 \qquad
 c_F=\prod_p\left(1-\frac{\nu_F(p)}{p^2}\right)>0.
\end{equation}
This $c_F$ is a density in the two affine parameters; it is \emph{not}
the Bateman--Horn singular series.

\textit{Proof.} Fix $z\ge2$. For each prime $p\le z$, exactly
$\nu_F(p)$ of the $p^2$ residue pairs $(a,b)\bmod p$ satisfy both
$p\mid a$ and $p\mid F(b)$. By the Chinese remainder theorem, the number
of pairs avoiding these simultaneous divisibilities at all such primes is
\begin{equation}\label{ra:54:finitecrt}
 N_{F,z}(T)=c_{F,z}T^2+O_{F,z}(T+1),\qquad
 c_{F,z}=\prod_{p\le z}\left(1-\frac{\nu_F(p)}{p^2}\right).
\end{equation}
Any pair counted by $N_{F,z}$ but not by $N_F$ has a prime $p>z$
dividing $a$, and necessarily $p\le T$. A union bound gives
\begin{align}
 0\le N_{F,z}(T)-N_F(T)
 &\le\sum_{z<p\le T}\left\lfloor\frac{T}{p}\right\rfloor
       \nu_F(p)\left(\frac{T}{p}+1\right)\notag\\
 &\le\frac{\Delta T^2}{z-1}
       +\Delta T(1+\log T).\label{ra:54:tail}
\end{align}
Only the elementary bounds on the sums of $n^{-2}$ and $n^{-1}$ were
used in the second line. Every factor defining $c_F$ is positive, and
$\sum_p\nu_F(p)/p^2<\infty$, so the infinite product converges to a
positive number. First divide \eqref{ra:54:finitecrt}--\eqref{ra:54:tail}
by $T^2$ and take upper and lower limits as $T\to\infty$ with $z$ fixed.
Then let $z\to\infty$. The error $\Delta/(z-1)$ disappears and both
bounds converge to $c_F$. This justifies the order of limits and proves
\eqref{ra:54:density}. \hfill$\square$

\textbf{Lemma 54.2 (preservation and injectivity).} Whenever
$\gcd(a,F(b))=1$ and $a\ge1$, the tuple
\[
 \mathbf f_{a,b}(x)=\bigl(f_1(ax+b),\ldots,f_r(ax+b)\bigr)
\]
is admissible, consists of primitive irreducible integer polynomials,
and retains positive leading coefficients and the same degree vector.
Different positive $a$ and different $b$ give different ordered tuples.

\textit{Proof.} If $p\nmid a$, the substitution $n\mapsto an+b$
permutes $\F_p$, so an allowed residue for $F$ gives one for
$F(ax+b)$. If $p\mid a$, the gcd hypothesis gives $F(b)\not\equiv0$
modulo $p$, and every residue is allowed. This proves admissibility.
Composition with an invertible affine map over $\Q$ preserves
irreducibility over $\Q$. If a transformed polynomial had a nontrivial
integer content, a prime dividing that content would divide the product
at every integer, contrary to admissibility; hence it is primitive.
The remaining sign and degree assertions follow from $a>0$.

For injectivity choose a member of degree $D$ and write its leading
terms as $c_Dx^D+c_{D-1}x^{D-1}$. In its transform these two coefficients
are
\[
 c_Da^D,\qquad a^{D-1}(D c_D b+c_{D-1}).
\]
The first determines the positive integer $a$, and the second then
determines $b$. This also works for $D=1$, when the second coefficient
is the constant coefficient. \hfill$\square$

There is a constant $C_{\mathbf f}\ge1$ such that, for $1\le a,b\le T$,
\begin{equation}\label{ra:54:height}
 H(\mathbf f_{a,b})\le C_{\mathbf f}T^D,
\end{equation}
where height is the maximum absolute coefficient. Indeed, the coefficient
of $x^k$ in $f_i(ax+b)$ is
$\sum_{j=k}^{d_i}c_{i,j}\binom jk a^k b^{j-k}$, each summand being
bounded by a fixed constant times $T^{d_i}$. Summing finitely many
coefficients proves the bound. No claim about equidistribution of these
coefficients is made.

\paragraph{Attempt: amplify a putative counterexample.}
Fix the degree vector $\mathbf d=(d_1,\ldots,d_r)$, and let
$E_{\mathbf d}(H)$ count admissible ordered tuples of distinct primitive
irreducible integer polynomials of these degrees, with positive leading
coefficients and height at most $H$, for which there is \emph{no}
positive integer giving simultaneous prime values.

\textbf{Proposition 54.3 (bad-orbit amplification).} If Hypothesis H
fails for one tuple $\mathbf f$ of degrees $\mathbf d$, then there are
constants $c_{\mathbf f}'>0$ and $H_{\mathbf f}$ such that
\begin{equation}\label{ra:54:amplification}
 E_{\mathbf d}(H)\ge c_{\mathbf f}'H^{2/D}
 \qquad(H\ge H_{\mathbf f}).
\end{equation}

\textit{Proof.} Failure of the original infinitude conclusion means
that some $B$ is larger than every positive integer at which the tuple
is simultaneously prime. For $b>B$, $a\ge1$ and $n\ge1$, we have
$an+b>B$. Thus every $\mathbf f_{a,b}$ has no simultaneous prime value
at a positive integer. Among the $(a,b)$ with $1\le a,b\le T$, discarding
$b\le B$ removes at most $BT=o(T^2)$ pairs. Lemma~54.1 leaves
$(c_F+o(1))T^2$ admissible pairs, Lemma~54.2 gives that many distinct
bad tuples, and \eqref{ra:54:height} puts them at height at most
$C_{\mathbf f}T^D$. Take
$T=\lfloor(H/C_{\mathbf f})^{1/D}\rfloor$ and decrease a positive
constant to absorb the integer part and the $o(1)$ term. This proves
\eqref{ra:54:amplification} for every sufficiently large $H$.
\hfill$\square$

A rigorous sufficient improvement of coefficient-averaging is therefore
\begin{equation}\label{ra:54:missingestimate}
 \boxed{\quad E_{\mathbf d}(H)=o(H^{2/D})
 \quad\text{for every fixed degree vector }\mathbf d.\quad}
\end{equation}
If this estimate were established, Proposition~54.3 would rule out every
fixed counterexample and prove the exact qualitative target. No
Bateman--Horn asymptotic is demanded of the resulting fixed tuple.

The threshold is substantially sharper than the one obtained using
translations alone, which would yield only $H^{1/D}$ bad tuples.
The gain comes from proving that a positive proportion of the two
parameters preserve admissibility, not from assuming that all dilations
preserve it. This is a proved notebook deduction; no priority claim is
made for this elementary amplification argument.

\paragraph{Computation: check the parameter mechanism, not the conjecture.}
For three admissible products $F(b)=b$, $b^2+1$, and $b(b+2)$,
exact enumeration gives the following numbers of good pairs:
\[
\begin{array}{c|rrr}
 T & F(b)=b & F(b)=b^2+1 & F(b)=b(b+2)\\\hline
 10&63&69&49\\
 100&6087&6729&4811\\
 1000&608383&671349&484132
\end{array}
\]
The following standard-library Python code reproduces the counts; it
uses integer gcd, not floating-point primality tests.
\begin{verbatim}
from math import gcd
for T in (10, 100, 1000):
    counts = []
    for F in (lambda b: b, lambda b: b*b+1,
              lambda b: b*(b+2)):
        counts.append(sum(gcd(a, F(b)) == 1
                          for a in range(1, T+1)
                          for b in range(1, T+1)))
    print(T, *counts)
\end{verbatim}
These finite checks support the parameter bookkeeping only. The density
and the amplification were proved above; the enumeration gives no
information about whether $n^2+1$ or $(n,n+2)$ has infinitely many prime
values.

\paragraph{Stress test: why the averaging theorem still does not transfer.}
The orbit just constructed has only order $H^{2/D}$ points in a much
larger coefficient space. Write $N=\sum_i(d_i+1)$ for the number of
coefficients in a full box. Even an upper allowance of the form
$H^N/(\log H)^A$, for every fixed $A$, is much larger than
$H^{2/D}$ whenever $N>2/D$. It can contain an entire bad affine orbit.
The logarithmic exceptional-set savings in the averaging literature
therefore do not supply \eqref{ra:54:missingestimate}; the allowed
coefficient cubes cannot simply be replaced by these curved orbits.
\eref{54}{2}

The exceptional corner is $r=1,D=1$, where $N=2=2/D$ and an appropriate
logarithmic saving could beat this obstruction. That is consistent with
the already solved single-linear-polynomial case, not evidence for the
higher-degree or multiple-linear cases. More generally, a theorem
excluding complete exceptional affine orbits might be sufficient without
a global bound as strong as \eqref{ra:54:missingestimate}, but such a
nonconcentration theorem is also unproved here. Calling the orbit
``random'' would not establish it.

A finite covering construction does not disprove H either. For any fixed
finite set $\mathcal P_0$ of primes, choose for each $p\in\mathcal P_0$
a residue $b_p$ with $F(b_p)\not\equiv0\pmod p$. Chinese remaindering
produces infinitely many integers avoiding all these prime divisors.
This proves only finite-level avoidance. It does not permit the exchange
of $\forall$ finite prime sets with one integer avoiding all primes, and
the surviving values can still be composite with larger prime factors.

\paragraph{Stress test: explicit audit of the cited solution claim.}
The transition to equation~(3.9) of Xiao's preprint treats
$|\mu(f_i(n))|$ as periodic along a residue class modulo a divisor $d$.
The congruence $d\mid F(n)$ is periodic, but squarefreeness is not.
Already for the admissible polynomial $f(x)=x$, the asserted right-hand
side becomes
\[
 \sum_{d\ge1}\mu(d)\log d\,|\mu(d)|\,
       \frac{z^d}{1-z^d}.
\]
Its coefficient at $z^4$ is $-\log2$, whereas the left-hand side
$H_f(z)=-\sum_{p\text{ prime}}\log p\,z^p$ has coefficient zero there.
Thus this specific identity is false, independently of any conjectural
prime asymptotic. \sref{54}{3}, \eref{54}{4}

There is a separate defect in the proposed final limiting justification
following its equation~(3.14): a function's being a polynomial does not
justify termwise limits in an unrelated infinite-series representation.
\eref{54}{4}
For an elementary diagnostic, define on $0\le z<1$
\[
 a_1(z)=1,\qquad a_n(z)=z^{n-1}-z^{n-2}\quad(n\ge2).
\]
The series is absolutely convergent for each $z<1$ and sums to the zero
polynomial, since its $N$th partial sum is $z^{N-1}$. Yet
\[
 \lim_{z\uparrow1}\sum_{n\ge1}a_n(z)=0,
 \qquad
 \sum_{n\ge1}\lim_{z\uparrow1}a_n(z)=1.
\]
Pointwise absolute convergence, even to a polynomial, supplies no
uniform tail estimate near the boundary. These tests reject the cited
argument, not Schinzel's hypothesis, and they also explain why no
analogous interchange was used in Lemma~54.1.

\paragraph{Dependencies.}
A successful proof for all tuples would imply Dickson's conjecture in
Section~\ref{62} by specializing to linear polynomials, including the
pair $(n,n+2)$. Conversely, the fixed-family Bateman--Horn asymptotic of
Section~\ref{83}, with its positive constant, would imply this qualitative
infinitude assertion after removing duplicates. The proposed
exceptional-set estimate is not derived from either unresolved statement.

\paragraph{Outcome.}
\textbf{Outcome: Conditional reduction.}

The affine parameter density, preservation of admissibility, injectivity
of the parameterization, height bound and $H^{2/D}$ counterexample
amplification were proved. They yield the explicit sufficient estimate
\eqref{ra:54:missingestimate}. The finite computation checks that
parameter mechanism, while the source audit identifies two independent
failures in one catalogue proof claim.

No exceptional-set estimate at the required strength was obtained, and
no counterexample family was found. Existing coefficient-average results
leave enough exceptional room for the entire affine orbit of a single
bad tuple. The remaining task is therefore not a routine passage from
``almost all'' to ``every'': it is a new quantitative estimate or a
structural theorem forbidding such exceptional orbits.
% END RESEARCH ADDITION ATTEMPT-54
\subsection{Sources}
\begin{itemize}\raggedright
\item[\textnormal{[S1]}]\hypertarget{source:54:1}{} Federico Scavia, Arithmetic Geometry and Algebraic Groups, Conjecture 3.1 (2026).
\url{https://www.lms.ac.uk/sites/default/files/inline-files/Scavia%20-%20Arithmetic%20Geometry%20and%20Algebraic%20Groups.pdf}.
{\small\textit{Catalogue formulation source.}}
\item[\textnormal{[S2]}]\hypertarget{source:54:2}{} On The Schinzel–Wójcik Problem Under Hypothesis H.
\url{https://www.cambridge.org/core/journals/mathematical-proceedings-of-the-cambridge-philosophical-society/article/abs/on-the-schinzelwojcik-problem-under-hypothesis-h/C75A7C818639ADA780DE2A47C8BCFA19}.
{\small\textit{Catalogue research/status reference; not a proof endorsement.}}
\item[\textnormal{[S3]}]\hypertarget{source:54:3}{} From Golomb to Schinzel.
\url{https://www.preprints.org/manuscript/202505.0651}.
{\small\textit{Catalogue research/status reference; not a proof endorsement.}}
% BEGIN RESEARCH ADDITION REFERENCES-54
\item[\textnormal{[E1]}]\hypertarget{extra:54:1}{} Alexei N.\ Skorobogatov and Efthymios Sofos, \emph{Schinzel Hypothesis on average and rational points}, Inventiones Mathematicae 231 (2023), 673--739; arXiv:2005.02998; DOI 10.1007/s00222-022-01153-6. Introduction and Theorems 1.1--1.2.
\url{https://arxiv.org/abs/2005.02998}.
{\small\textit{Added research source: Fixed-family formulation and coefficient-averaged existence of a simultaneous prime value; the quantifier limitation is retained. Consulted 24 September 2026.}}
\item[\textnormal{[E2]}]\hypertarget{extra:54:2}{} Tim Browning, Efthymios Sofos and Joni Ter\"av\"ainen, \emph{Bateman--Horn, polynomial Chowla and the Hasse principle with probability 1}, arXiv:2212.10373v2, 21 May 2026, Theorems 1.2 and 1.4 and the introduction.
\url{https://arxiv.org/abs/2212.10373v2}.
{\small\textit{Added research source: Averaged asymptotics and logarithmic exceptional-proportion savings in coefficient boxes; no result along every fixed affine orbit is assumed. Consulted 24 September 2026.}}
\item[\textnormal{[E3]}]\hypertarget{extra:54:3}{} John Friedlander and Henryk Iwaniec, \emph{Asymptotic sieve for primes}, Annals of Mathematics 148 (1998), 1041--1065; arXiv:math/9811186, Section 1.
\url{https://arxiv.org/abs/math/9811186}.
{\small\textit{Added research source: Parity obstruction and the need for additional analytic information beyond ordinary congruence counts. Consulted 24 September 2026.}}
\item[\textnormal{[E4]}]\hypertarget{extra:54:4}{} Huan Xiao, \emph{From Golomb to Schinzel}, Preprints.org manuscript 202505.0651, version 2, posted 9 June 2025, Section 3, especially the transition from (3.8) to (3.9) and the paragraph following (3.14).
\url{https://www.preprints.org/manuscript/202505.0651}.
{\small\textit{Added research source: Claimed solution, not endorsed. The attempt gives a coefficient-level counterexample to the displayed identity and an independent test of the limiting justification. Consulted 24 September 2026.}}
% END RESEARCH ADDITION REFERENCES-54
\end{itemize}

\end{document}
